\chapter*{Brief Summary}
The COMETES ASM Team Project 2026 focuses on the development and integration of a mechatronic satellite subsystem. The system under study consists of a at satellite (flatsat) to which additional components, namely an antenna and a radar module, are added. Based on these extensions, safety aspects must be analyzed and considered throughout the design process. The project is carried out in collaboration with students from the French partner university SUPMECA in Paris. It combines mechanical design, system integration, simulation, and safety analysis within an international team environment.\\
\\
The scope was first defined as the development of a deployable antenna and radar system for a flat satellite. Then requirements were clarified with a stakeholder and a context diagram for the geometry, safety and environmental conditions.\\
\\
To develop the model, it has been decided that a CAD model should be designed to be then imported in a multi-body model for simulations in Matlab.
In parallel, components for a self-locking system have be selected to develop a self-locking transmission which actuates the movement of the multi-body antenna regarding the requirements of safety.\\
\\
Throughout the entire project, multiple online meetings  with Supm\'{e}ca students took place in order to update requirements, progress of both teams,share ideas and compare results. The parisian team developed a lipstick mechanism to deploy the antenna while it has been chosen in Germany to implement ropes along the ribs of the antenna in order to enable their retractation additionally to their deployment.\\
\\
At the end of the given time period before the common internatonal presentation, the biggest challenge was to join the multibody model and the control architecture. Seperated, they give satisfactory results. However, when they are joined, the model complexity level increases so much that the simulation gets very slow. Moreover, some interferences problems must be solved.\\ \\
Finally, this project allowed to develop an advanced physical model. If other students want to pursue its development in the future, their next steps would be a further analysis of our model in different situations as well as manufacturing possibilities checks. Then manufacturing different components. Implement physical tests on sub-components. Assemble the system. Implement tests on the assembled system and then validate it.

